\subsection{Exact common powers and the joint derivative budget}

The finite and radial preparations label the same outgoing orbits in
different ways: by $p$ on $Y=0$ and by the horizontal phase.  We
identify these labels through a regular connector before applying an
upper slow passage.  This gives a common input for the two reflected
lower powers.  Their exact exponents agree for the same field, so
their singular power factors cancel in the logarithmic derivative
ratio.

\begin{lemma}[A common actual input]\label{rw:transitions:common-input}
Fix a compact coefficient box and a fixed annulus
$s_1\le\sigma\le2s_1$, where $\sigma=1/p$ and $p$ labels the
outgoing intersection with $Y=0$. Suppose this annulus and a slightly
larger one lie in the common low-layer connector to the selected
horizontal section, before every upper slow block under consideration.
Let $t$ be the reflected lower phase on that horizontal section.
For each fixed finite order $N$,
\[
 t=t_\epsilon(\sigma),\qquad
 \|t_\epsilon-\sigma\|_{C^N_{\sigma,\mathbf b},E_\epsilon}
 \le C_N\epsilon(1+|\log\epsilon|)^{M_N}=o(1).
\]
There is consequently a common coordinate on these outgoing orbits
which agrees with $t$ near the radial end and with $1/p$ on the
finite side. On the overlap its selected jets and inverse jets are
bounded and its derivative is positive. The choice is uniform in
subsequent slow gaps, provided the stated connector remains in the
fixed common domain.
\end{lemma}
\begin{proof}
Use the reflected version of the lower phase and reverse its time.
On $Y=0$ its phase is $\sigma+O_{C^N}(\epsilon)$ and its normal
input is $Z_e=\epsilon A$, where $A$ is a positive bounded unit on
this fixed annulus. Its backward normal velocity is $Z\beta_I$.
Follow it to the fixed positive horizontal section, retaining the
slightly larger annulus. In the normal logarithm
$r=\log(Z/Z_e)$, the phase equation is
\[
 \frac{\dd x}{\dd r}
   =-\frac{\epsilon g_I(x)}{\beta_I(x,Z_e e^r)}.
\]
The terminal value of $r$ is $\log(1/\epsilon)+O(1)$, since the
horizontal section has a positive normal height. Its transversality
has a fixed margin. Throughout this connector $Z_e e^r$ stays in
the fixed normal tube. The phase moves by
$O(\epsilon(1+|\log\epsilon|))$, so the larger annulus indeed
contains it when the scale is small.

For the mixed rows, the highest phase variation has coefficient
$O(\epsilon)$ and its forcing is $\epsilon$ times bounded lower
variations. The selected derivatives of $\log Z_e$ are bounded and
$Z_e e^r$ stays in the tube. The scalar variational induction and
Gronwall on this $O(1+|\log\epsilon|)$ interval therefore give the
claimed $C_N\epsilon(1+|\log\epsilon|)^{M_N}$ bound. The fixed
terminal transversality controls its endpoint rows; no slow denominator
is inverted.

On the fixed overlap choose a smooth cutoff $\chi$, equal to one
at its radial side and zero at its finite side, and use
\[
 \sigma+\chi(\sigma)(t_\epsilon(\sigma)-\sigma).
\]
Its derivative is $1+o(1)$, and finite-order inverse differentiation
gives the inverse bounds. Continue it by $t$ toward the radial end
and by $\sigma$ toward the finite side. These are coordinates on
the same actual outgoing orbits, obtained entirely within the regular
connector. Subsequent coefficient paths and passing conditions are
chosen inside these common charts.
\end{proof}

We now check the exponent equality needed for this cancellation.
In the horizontal chart used above, the weak and normal
rates at its origin are $\epsilon\bar B$ and $-(1-\epsilon\bar B)$.
Reflection fixes both $\epsilon$ and $\bar B$. Thus, on a compact set
$\bar B\ge B_*>0$, the two horizontal passages have the same exact
exponent
\begin{equation}\label{rw:transitions:equal-exponents}
 \mu=\frac{1-\epsilon\bar B}{\epsilon\bar B},\qquad
 h=\epsilon\mu=\frac{1-\epsilon\bar B}{\bar B}>h_*>0.
\end{equation}
The weak boundary $\bar B=0$ is treated by the unnormalized slow
factor of Proposition~\ref{rw:lower-preparation}.

\begin{lemma}[Cancellation before differentiation]\label{rw:joint-jet-budget}
Let $D_\pm(s)=K_\pm s^{h/\epsilon}e^{P_\pm(s)/\epsilon}$,
$K_\pm>0$, on $s>0$, with the same $h$ in both maps. Let
$f(s)=s k(s)$, $k>0$, and suppose the following three factors have
positive lower bounds:
\[
 1+\frac{s k'}k,\qquad h+sP_+'(s),\qquad
 h+f(s)P_-'(f(s)).
\]
Assume bounded total order $N+1$ for $k,P_\pm$ and order $N$ for
$h,1/k$, on a common input and coefficient neighborhood. Then
\begin{equation}\label{rw:transitions:common-power}
\begin{aligned}
 \epsilon\log\frac{(D_-\circ f)'(s)}{D_+'(s)}
 ={}&C(\mathbf b,\epsilon)+h\log k(s)
       +P_-(f(s))-P_+(s)\\
 &+\epsilon\log\left(1+\frac{s k'(s)}{k(s)}\right)
  +\epsilon\log\frac{h+f(s)P_-'(f(s))}{h+sP_+'(s)},\\
 C(\mathbf b,\epsilon)={}&\epsilon\log(K_-/K_+).
\end{aligned}
\end{equation}
After subtracting $C$, the right side has bounded total order $N$.
No boundedness assertion for $C$ follows from this cancellation.
\end{lemma}
\begin{proof}
Differentiate both amplitudes before cancelling factors. The power part
of the ratio is
\[
 (f/s)^{h/\epsilon-1}f'
       =k^{h/\epsilon}\left(1+\frac{s k'}k\right).
\]
The remaining factors are $h+fP_-'(f)$ and $h+sP_+'(s)$; logarithms
give the identity. Their positive margins bound the reciprocals and
logarithms. Each differentiated primitive costs one order, so the
total-order chain rule gives the stated budget, including coefficient
dependence of $f$ and Euler scale derivatives.
\end{proof}

The identity makes the derivative requirement explicit.  Primitive
total order five leaves total order four for the nonconstant
$R=K-K(0)$.  For $H=\partial_sR$, the orders $s^3\mathbf b^0$ and
$s^2\mathbf b^1$ give two input derivatives and coefficient Lipschitz
control of its first derivative. A Hadamard column additionally uses
$s^3\mathbf b^1$ and $s^2\mathbf b^2$, still total order four.
To use that column, one must also establish its center identity and
the common domain along the whole coefficient segment; these are
proved separately.

For example $k=\alpha e^\psi$, with an input-independent $\alpha>0$
that may be small, gives
\[
 h\log k=h\log\alpha+h\psi,\qquad
 1+s k'/k=1+s\psi'.
\]
Group $h\log\alpha$ with $C$; the remaining compositions use
bounded jets of $f,\psi$.  Derivatives of the coefficient-only
constant require separate action bounds.  Exact exponent equality is
essential: an error $\delta\mu$ leaves the unbounded term
$\epsilon\delta\mu\log f(s)$ at $s=0$.

The constructions must therefore reserve the derivative orders before
their domains are reduced.  The following allocation records the
orders used later.
\begin{center}\small
\begin{tabularx}{\linewidth}{@{}lXX@{}}
\toprule
Use & Primitive order and operations & Result used later\\
\midrule
Three lower connections & total order at least ten, including
six input, two coefficient and two Euler orders & bounded relative
units and their compositions on the stated common domains\\
Separated joint columns & total order five; one derivative in
\eqref{rw:transitions:common-power} & total order four for the
nonconstant comparison, including $s^3\mathbf b^1$ and $s^2\mathbf b^2$\\
Contact comparison & total order eight if three input and three
coefficient derivatives of its numerator are retained & at most two
orders spent on logarithmic rates and connection curvature\\
Weak or cubic rate & total order eight; two rate derivatives and
one later Hadamard division & three input and two coefficient orders
of the divided numerator, once that division is justified\\
Upper saddle & primitive total order six for the relative logarithm;
total order eight for its curvature & four Euler and two coefficient
orders of the respective object; curvature consumes two extra Euler
orders\\
\bottomrule
\end{tabularx}
\end{center}
The exact long-scale maps and their ordinary terminal-height derivatives
are constructed separately.
